\documentclass[11pt,a4paper]{article}
\usepackage[T1]{fontenc}
\usepackage[utf8]{inputenc}
\usepackage{amsmath,amssymb,xcolor}
\usepackage[margin=25mm]{geometry}
\usepackage[colorlinks=true,urlcolor=blue]{hyperref}
\definecolor{nexorablue}{HTML}{2459B5}
\setlength{\parindent}{0pt}
\setlength{\parskip}{7pt}
\setlength{\emergencystretch}{2em}
\renewcommand{\familydefault}{\sfdefault}
\begin{document}

{\small\bfseries\color{nexorablue} NEXORA CPGE}\par
{\small MPSI / NEX-LOG-01 / Corrigé}\par
{\LARGE\bfseries Décrire le rejet par les deux tests}\par
\vspace{0.25em}\hrule\vspace{0.6em}

\section*{Négation}
Le rejet signifie que $n$ n’est pas multiple de $3$, ou que $n\le-4$, ou que $n>8$. Les bornes viennent de la négation de $-4<n\le8$.

\section*{Témoins}
Le nombre $7$ est dans l’intervalle mais n’est pas multiple de $3$. Le nombre $9$ est multiple de $3$ mais est hors de l’intervalle. Ces deux entiers isolent les deux causes.


\end{document}
